\section{Checks against misreadings}\label{sec:integrated-stress}



In this section a misreading is \emph{excluded} when it is not licensed by the
results of
Sections~\ref{sec:fatcd-domain}--\ref{sec:decomposition-exhaustion}.
The section is a checklist. Each of its six remarks lists misreadings of
the results of Sections~\ref{sec:lower-bounds}--\ref{sec:decomposition-exhaustion},
for example asserting exact six without the lower bounds, or reading a
role channel as an active role, and names the earlier result that rules
each misreading out; ``stress'' means a check against such a list, and the
chain \emph{satisfies} a stress family when the result named for each
misreading in its list rules that misreading out.
Proposition~\ref{thm:integrated-stress} checks that the earlier results
block the misreadings listed in the six remarks and combine into the terminal statement.
The section adds no new mathematics; it collects in one place what the
results do not permit.

\begin{center}
\small
\begin{tabularx}{\linewidth}{@{}>{\raggedright\arraybackslash}p{0.34\linewidth}>{\raggedright\arraybackslash}X@{}}
\toprule
Misreadings of & Blocked by \\
\midrule
dependency (Remark~\ref{prop:dependency-stress}) & the \FATCD{} domain and the order of Theorems~\ref{thm:six-role-lower-bounds}, \ref{thm:upper-bound-classification}, \ref{thm:decomposition-existence} and~\ref{thm:six-channel-exactness} \\
circularity (Remark~\ref{prop:circularity-stress}) & Proposition~\ref{prop:non-circularity} and the direct proof of Theorem~\ref{thm:decomposition-existence} \\
channel versus activation (Remark~\ref{prop:channel-activation-stress}) & Proposition~\ref{prop:channel-vs-activation} and Definition~\ref{def:channel-status} \\
role alignment (Remark~\ref{prop:role-alignment-stress}) & Proposition~\ref{prop:alignment-rules} \\
upper-bound residual (Remark~\ref{prop:residual-stress}) & Propositions~\ref{prop:probability-closure}--\ref{prop:agency-closure} and Theorem~\ref{thm:upper-bound-classification} \\
overclaim (Remark~\ref{prop:overclaim-stress}) & Remark~\ref{rem:atlas-global-nonclaims} and Propositions~\ref{prop:nonclaim-closure} and~\ref{prop:non-uniqueness} \\
\bottomrule
\end{tabularx}
\end{center}

\subsection{Dependency stress}\label{subsec:dependency-stress}

\begin{remark}[Dependency stress]\label{prop:dependency-stress}
The scoped exact-six target \(\ScopedExactSix{\FATCD_{\mathrm{RH}}}\)
depends on the \FATCD{} domain of Section~\ref{sec:fatcd-domain}, the
upper-bound classification of Theorem~\ref{thm:upper-bound-classification},
and the decomposition and exhaustion theorems of
Section~\ref{sec:decomposition-exhaustion}; the last clause of
Theorem~\ref{thm:scoped-exact-six} also depends on the lower-bound theorem of
Section~\ref{sec:lower-bounds}. Accordingly, the following are excluded:
\begin{itemize}[topsep=2pt,itemsep=1pt]
  \item asserting the target without the \FATCD{} domain;
  \item asserting the target without the upper-bound classification;
  \item asserting the target without the decomposition and exhaustion
    theorems;
  \item asserting that every channel is active somewhere in
    \(\FATCD_{\mathrm{RH}}\) without the lower-bound theorem.
\end{itemize}
The integrated stress check of Proposition~\ref{thm:integrated-stress} is a
separate consistency check; the proof of Theorem~\ref{thm:scoped-exact-six}
does not use it.
\end{remark}

\subsection{Circularity stress}\label{subsec:circularity-stress}

\begin{remark}[Circularity stress]\label{prop:circularity-stress}
Three circularity patterns are excluded under the theorem chain:
\begin{itemize}[topsep=2pt,itemsep=1pt]
  \item redefining \FATCD{} so that it already contains exactly six role
    slots by construction;
  \item declaring the absence of a seventh role by a membership condition
    of \FATCD{} rather than by an explicitly stated hypothesis;
  \item justifying decomposition existence by assuming scoped
    exact-six.
\end{itemize}
The non-circularity schema of Proposition~\ref{prop:non-circularity} and
the direct proof of Theorem~\ref{thm:decomposition-existence}
block each pattern.
\end{remark}

\subsection{Channel-versus-activation stress}\label{subsec:channel-activation-stress}

\begin{remark}[Channel-versus-activation stress]\label{prop:channel-activation-stress}
The following conflations are excluded:
\begin{itemize}[topsep=2pt,itemsep=1pt]
  \item requiring that every \(D \in \FATCD{}\) activates all six role
    channels;
  \item treating an inactive channel of status \(\belowthr\) as an
    active witness for its role;
  \item treating an inactive channel of status \(\collapsedbc\) as an
    active witness;
  \item treating an inactive channel of status \(\absentrec\) as an
    active witness.
\end{itemize}
Proposition~\ref{prop:channel-vs-activation} enforces each distinction.
The same inactive-status check also applies to
\(\outsidescope\).
\end{remark}

\subsection{Role-alignment stress}\label{subsec:role-alignment-stress}

\begin{remark}[Role-alignment stress]\label{prop:role-alignment-stress}
The three forbidden assignments of
Remark~\ref{rem:three-repaired-alignments} are excluded: generic
transition data does not project to \Pthree{}, generic path or
derivation data does not project to \Pfive{}, and generic semantic or
interpretation data does not project to \Pone{}. The alignment rules
of Proposition~\ref{prop:alignment-rules} enforce each prohibition.
\end{remark}

\subsection{Upper-bound residual stress}\label{subsec:upper-bound-residual-stress}

\begin{remark}[Upper-bound residual stress]\label{prop:residual-stress}
Under Theorem~\ref{thm:upper-bound-classification} and within the covered
scope, the following misreadings are excluded:
\begin{itemize}[topsep=2pt,itemsep=1pt]
  \item probability or uncertainty features of the recognized forms of
    Proposition~\ref{prop:probability-closure} reopening as unresolved
    threats;
  \item causality features of the recognized forms of
    Proposition~\ref{prop:causality-closure} reopening as unresolved threats;
  \item agency features of the recognized forms of
    Proposition~\ref{prop:agency-closure} reopening as unresolved threats;
  \item any residual feature in the covered scope of a description satisfying the
    recognition hypothesis remaining unclassified under
    Definition~\ref{def:residual-scheme};
  \item any residual feature in the covered scope of a description satisfying
    the recognition hypothesis passing the recorded seventh-role screen
    (without that hypothesis such a feature exists: the relation record \(r\) of
    \(D_{\mathrm{scr}}\), Remark~\ref{rem:recognition-open}).
\end{itemize}
Propositions~\ref{prop:probability-closure}--\ref{prop:agency-closure}
settle the three threat closures, and
Theorem~\ref{thm:upper-bound-classification}, under the recognition
hypothesis, delivers the emptiness of categories (x) and (xi).
\end{remark}

\subsection{Overclaim stress}\label{subsec:overclaim-stress}

\begin{remark}[Overclaim stress]\label{prop:overclaim-stress}
The following restatements of the scoped exact-six target are
excluded:
\begin{itemize}[topsep=2pt,itemsep=1pt]
  \item restatement as a full six-primitives algebra;
  \item restatement as a global primitive-independence theorem;
  \item restatement as unique or canonical decomposition;
  \item restatement as empirical realization;
  \item restatement as a claim about all emergence phenomena, or any
    domain broader than \FATCD{}.
\end{itemize}
Remark~\ref{rem:atlas-global-nonclaims},
Proposition~\ref{prop:nonclaim-closure}, and
Proposition~\ref{prop:non-uniqueness} carry the corresponding nonclaims
forward through the theorem chain. In particular, nonclaim closure
survives the exact-six assembly.
\end{remark}

\subsection{The integrated stress certification}\label{subsec:integrated-stress-theorem}

\begin{proposition}[Integrated stress certification]\label{thm:integrated-stress}
Clauses (a)--(d) are statements about descriptions; clause (e) records the
dependency structure of the proof of Theorem~\ref{thm:scoped-exact-six}.
The following hold:
(a) every \(D\in\FATCD{}\), including \(D_{\mathrm{norm}}\notin\FATCD_{\mathrm{RH}}\),
admits a well-formed decomposition record, and every well-formed record of such
\(D\) satisfies clauses (i)--(ii) of Theorem~\ref{thm:six-channel-exactness};
(b) for each \(k\in\{0,\dots,6\}\) some description in \(\FATCD_{\mathrm{RH}}\) has exactly
\(k\) active channels (Proposition~\ref{prop:channel-vs-activation});
(c) some description has two distinct well-formed decomposition records
(Proposition~\ref{prop:non-uniqueness});
(d) the description \(D_{\mathrm{norm}}\) of Remark~\ref{rem:recognition-open}
lies in \FATCD{} but not in \(\FATCD_{\mathrm{RH}}\); and
(e) Theorem~\ref{thm:scoped-exact-six} follows from
Theorems~\ref{thm:six-role-lower-bounds},
\ref{thm:upper-bound-classification}, \ref{thm:decomposition-existence}
and~\ref{thm:six-channel-exactness}, Propositions~\ref{prop:probability-closure}--\ref{prop:agency-closure},
and the inventory of the records of \(D_1,\ldots,D_6\) given in its proof.
\end{proposition}

\begin{proof}
We verify that each of the six stress families is discharged by the earlier
result it rests on, so that none of the listed misreadings can arise in
the combined chain.

\emph{Dependency.} The components named in
Remark~\ref{prop:dependency-stress} are present and ordered: the
\FATCD{} domain is fixed in Section~\ref{sec:fatcd-domain}, the
lower-bound theorem in Section~\ref{sec:lower-bounds}
(Theorem~\ref{thm:six-role-lower-bounds}), the upper-bound
classification in Section~\ref{sec:upper-bound-classification}
(Theorem~\ref{thm:upper-bound-classification}), and the
decomposition and exhaustion theorems in
Section~\ref{sec:decomposition-exhaustion}
(Theorems~\ref{thm:decomposition-existence}
and~\ref{thm:six-channel-exactness}). Theorem~\ref{thm:scoped-exact-six}
uses Theorems~\ref{thm:decomposition-existence}
and~\ref{thm:six-channel-exactness}; clause (iii) of
Theorem~\ref{thm:six-channel-exactness} uses
Theorem~\ref{thm:upper-bound-classification}; and the last clause of
Theorem~\ref{thm:scoped-exact-six} uses
Theorem~\ref{thm:six-role-lower-bounds}. Theorem~\ref{thm:upper-bound-classification} does not use
Theorem~\ref{thm:six-role-lower-bounds}. The existence argument of
Theorem~\ref{thm:decomposition-existence} does not use it either; that proof
cites it only for the separate non-vacuity observation at its end. So the target cannot be asserted
while skipping one of these uses, and no link consumes its own output.

\emph{Circularity.} The three circularity patterns of
Remark~\ref{prop:circularity-stress} are blocked by the
non-circularity schema of Proposition~\ref{prop:non-circularity}, which
defines \FATCD{} without reference to a six-slot count, and by the
direct proof of Theorem~\ref{thm:decomposition-existence}, which
produces a decomposition without assuming scoped exact-six. The absence of
labels (x) and (xi) is not a membership condition of \FATCD{}: it follows,
through the label criteria of Definition~\ref{def:residual-scheme}, from the
recognition hypothesis, which is assumed and fails for the description of
Remark~\ref{rem:recognition-open}.

\emph{Channel versus activation.} The conflations of
Remark~\ref{prop:channel-activation-stress} are excluded by
Proposition~\ref{prop:channel-vs-activation} together with the
five-status taxonomy of Definition~\ref{def:channel-status}: a channel
carrying status \(\belowthr\), \(\collapsedbc\), \(\absentrec\), or
\(\outsidescope\) is an inactive channel and is never read as an active
projection, so no description is required to activate all six channels.

\emph{Role alignment.} The three forbidden assignments of
Remark~\ref{prop:role-alignment-stress} are ruled out by the alignment
rules of Proposition~\ref{prop:alignment-rules}, which refuse generic
transition data as \Pthree{}, generic path or derivation data as
\Pfive{}, and generic semantic or interpretation data as \Pone{}, in
accordance with the alignments of
Remark~\ref{rem:three-repaired-alignments}.

\emph{Upper-bound residual.} The misreadings of
Remark~\ref{prop:residual-stress} are closed within the covered scope by
Propositions~\ref{prop:probability-closure}--\ref{prop:agency-closure},
which settle the probability, causality, and agency threat closures, and
by Theorem~\ref{thm:upper-bound-classification}, which --- under the
recognition hypothesis (Definition~\ref{def:recognition-hypothesis}) ---
classifies every covered residual feature and delivers the emptiness of
categories (x) and (xi); hence, under that hypothesis, no covered residual
reopens as an unresolved threat or passes the recorded seventh-role screen, while the
three principal threat classes are closed without that hypothesis, in the
recognized forms of
Propositions~\ref{prop:probability-closure}--\ref{prop:agency-closure}.

\emph{Overclaim.} The excluded restatements of
Remark~\ref{prop:overclaim-stress} are blocked because the nonclaims of
Remark~\ref{rem:atlas-global-nonclaims} are carried forward by
Proposition~\ref{prop:nonclaim-closure}, and non-uniqueness of the
decomposition is recorded by Proposition~\ref{prop:non-uniqueness}; thus
nonclaim closure survives the exact-six assembly, and no full-algebra,
global-independence, unique-decomposition, empirical-realization, or
broadened-domain reading is admitted.

Thus each misreading listed in the stress remarks is excluded by the result
cited for it; dependencies cohere, the channel-versus-activation distinction
holds and nonclaim closure holds, so the results combine into the terminal
statement of
Section~\ref{sec:scoped-exact-six}.

For the concrete clauses: (a) holds because Theorem~\ref{thm:decomposition-existence}
gives a well-formed decomposition record for every \(D\in\FATCD{}\), and its
existence construction and clauses (i)--(ii) of
Theorem~\ref{thm:six-channel-exactness} use only
Definitions~\ref{def:channel-status} and~\ref{def:decomposition-record} and the
label criteria after Definition~\ref{def:residual-scheme}, and clause (iii)
alone cites Theorem~\ref{thm:upper-bound-classification}; (b) is
Proposition~\ref{prop:channel-vs-activation}; (c) is
Proposition~\ref{prop:non-uniqueness}; (d) is shown in
Remark~\ref{rem:recognition-open}; and (e) holds because the proof of
Theorem~\ref{thm:scoped-exact-six} in Section~\ref{sec:scoped-exact-six} uses exactly
the results listed in (e) and not this proposition.
\end{proof}

In the sense of the checklist of this section, the results named in its table
rule out each listed misreading: dependency, circularity,
channel-versus-activation, role-alignment, upper-bound residual and overclaim
misreadings.
